\documentclass{article}
\usepackage{iclr2027_conference,times}
\input{math_commands.tex}
\usepackage{microtype}
\usepackage{hyperref}
\usepackage{url}
\usepackage{booktabs}
\usepackage{graphicx}
\usepackage{multirow}

\title{MolProgram: Compositional Property Programs for Molecular Construction\\
and Editing with Group-Relative Reinforcement Learning}

% Keep anonymous for ICLR 2027 submission.
\author{Anonymous authors}

% \iclrfinalcopy
\begin{document}
\maketitle

\begin{abstract}
Goal-directed molecular design is central to drug discovery and materials
development, yet practical requests combine a variable number of numerical
objectives and may ask either for a new molecule or for a source-preserving
edit. Existing systems typically encode these requests as flat property vectors
or free-form instructions, optimize de novo generation and editing in isolation,
and obscure policy quality through property-selected best-of-many evaluation.
We introduce \emph{MolProgram}, a molecular design framework that represents
each request as a variable-length program of property identities, targets,
tolerances, and directions. MolProgram shares this request language, its
programmatic verifier, and a group-relative reinforcement objective across two
state-appropriate transition operators: a direct SMILES constructor for an
empty initial state and a source-anchored graph-event kernel for a populated
state. This separation is determined by the boundary condition rather than by
a learned task router; it preserves a common optimization problem without
forcing molecular construction and local editing into the same output geometry.
The construction branch combines frozen LLM-derived instruction features with
explicit numerical program tokens and a property-count curriculum, while both
branches receive group-relative molecular feedback without retrieval from an
external molecule library or property-aware inference-time reranking. On direct
de novo generation, group-relative post-training improves Pass@8 and Pass@20
over supervised fine-tuning by 6.35 and 7.64 percentage points on
two-to-seven-property programs, and by 3.90 and 6.50 points on
out-of-distribution programs. On MolEdit-Instruct, group-relative post-training
raises the source-anchored editor's strict candidate-level accuracy from
43.35\% to 46.12\%, exceeding MolEditRL's reported 45.0\% macro average under
the same strict metric, while relaxed accuracy remains lower (61.90\% versus
72.7\%). These results show that compositional property programs can unify the
request and learning interface across molecular construction and editing while
retaining transition operators matched to their distinct chemical constraints.
\end{abstract}

\section{Introduction}
Goal-directed molecular design is widely used in drug discovery and materials
development to create or modify molecules toward desired physicochemical and
biological profiles. It is commonly formulated as a conditional molecular
transition problem: given a property program
$P=\{(j,y_j,\tau_j,d_j)\}_{j=1}^{m}$ and an initial molecular graph $G_0$,
learn a policy $\pi_\theta$ whose terminal graph satisfies all clauses in $P$.
Each clause specifies a property identity $j$, target $y_j$, tolerance
$\tau_j$, and optional direction $d_j$; $G_0=\varnothing$ defines de novo
generation, whereas $G_0=G_{\mathrm{src}}$ defines source-conditioned editing.
A general solution iteratively samples a molecular action
$a_t\sim\pi_\theta(a_t\mid P,G_t,a_{<t})$ from the state-dependent legal
support $\mathcal A(G_t)$ and applies an interpreter
$G_{t+1}=\mathcal I(G_t,a_t)$ until a stop action produces the candidate
molecule.

Existing approaches instantiate this iteration through several broad
paradigms rather than isolated model families. Sequence generators decode
molecular strings token by token and provide a compact route to conditional
generation \citep{bagal2022molgpt}. Graph and image diffusion models iteratively
denoise structural states and can impose numerical or categorical guidance
throughout construction \citep{liu2024graphdit,wang2025sketchmol}.
Arbitrary-subset conditional models use property masking, guidance, or
self-criticism to operate on different subsets of available objectives
\citep{jolicoeurmartineau2024stggplus}. Instruction-tuned molecular language
models use natural-language task descriptions to transfer across objectives
and molecular operations \citep{dey2025gellm3o,dey2025controllable}. In
parallel, verifier-guided sampling and reinforcement learning use molecular
predictors or programmatic property functions to improve outcomes that are not
fully captured by supervised likelihood
\citep{wang2025sketchmol,gao2026cmoral,shao2024deepseekmath}.

These paradigms all approximate a sequence of condition-dependent molecular
transitions; however, their interfaces make compositional and operational
generalization difficult to identify. A flat property vector does not expose
the internal structure or length of a variable objective set, making
high-order composition difficult to distinguish from interpolation among
individual property correlations. A free-form instruction provides semantic
flexibility but is imprecise about numerical targets, tolerances, and reversed
or extreme constraints. Direct string generation is efficient for constructing
whole molecules but fragile under source-preservation constraints, whereas
graph events maintain chemical locality but impose an unnecessarily long
horizon when constructing from an empty state. Treating the two settings as
unrelated discards their common objective language; forcing them through one
decoder or action space instead creates a transition-geometry conflict.
Best-of-many evaluation can hide either limitation by increasing success
through search or property-aware selection rather than by increasing the policy
probability of a strict solution \citep{gao2022sampleefficiency}.

Our key observation is that the reusable object is the property program and its
verifiable objective, not necessarily every parameter of the molecular
transition operator. The empty boundary requires global construction, whereas
a populated boundary requires local, source-anchored change. Sharing the
request representation and group-relative reward makes success mean the same
thing in both regimes; selecting the transition operator deterministically from
source occupancy preserves the error geometry appropriate to each. This is a
hybrid solver for one conditional transition problem, rather than a learned
router that decides which task is being requested.

Thus, we present \textbf{MolProgram}, which combines LLM-derived instruction
conditioning, explicit property-program tokens, state-appropriate molecular
transitions, and group-relative reinforcement learning. MolProgram first
normalizes every request into deterministic tokens for its active property
clauses. For an empty initial graph, a compact autoregressive Transformer emits
a complete SMILES string and is trained with a property-count curriculum. For
a populated graph, a validity-preserving event kernel proposes local graph
transitions anchored to the source. Both operators are optimized from groups of
raw candidates using the same programmatic property and source-similarity
verifier. Source occupancy chooses the operator before inference; neither an
LLM router nor an output-side property ranker selects among models or
candidates.

The main contributions of this work are as follows:
\begin{enumerate}
  \item We formulate de novo generation and molecular editing as one
  initial-state-conditioned transition problem and introduce MolProgram, which
  shares a variable-length property-program interface and verifiable objective
  across a direct molecular constructor and a source-anchored event kernel,
  without a learned task router, retrieval library, or property-aware
  inference-time selection.
  \item We develop property-count curriculum learning, verifier-based
  group-relative post-training, and a budget-transparent evaluation that
  separately reports raw success, Pass@$k$, validity, uniqueness, and
  source-preserving editing. The resulting evidence shows improved high-order
  and out-of-distribution de novo control and a strict candidate-level editing
  gain over MolEditRL, while retaining the relaxed-threshold deficit as an
  explicit limitation.
\end{enumerate}

The remainder of the paper is organized as follows. Section 2 reviews
conditional molecular generation, editing, and reinforcement post-training.
Section 3 formalizes compositional property programs and the unified molecular
transition problem. Section 4 presents MolProgram and its training objectives.
Section 5 reports de novo, editing, complexity-scaling, and
out-of-distribution results. Section 6 discusses limitations and broader
implications, and Section 7 concludes the paper.

\section{Related Work}
% To be written.

\section{Problem Formulation}
% To be written.

\section{Method}
% To be written.

% Results combine the frozen direct policy and source-anchored event-kernel evidence.
\section{Results}
\label{sec:results}

We organize the results around three questions. First, does group-relative
post-training improve compositional property control at practical candidate
budgets rather than only under large best-of-$k$ search? Second, does the gain
persist as property programs become longer and move out of distribution? Third,
can de novo generation and source-conditioned editing be expressed and learned
through one executable transition interface? Unless stated otherwise, all
candidate streams preserve generation order, Raw@1 is exactly the first model
draw, and property-reranked selection is never reported as one-shot generation.
All experiments use the frozen single seed specified before evaluation.

\subsection{Group-relative post-training improves low-budget de novo design}
\label{sec:results-budget}

Table~\ref{tab:budget-scaling} compares supervised fine-tuning (SFT) and
Group-RL on 6,000 two-to-seven-property programs and 1,000
out-of-distribution (OOD) programs. The improvement is already visible at
small candidate budgets. On 2p--7p, strict success on the first draw increases
from 7.02\% to 8.93\%; Pass@8 and Pass@20 increase by 6.35 and 7.64 percentage
points. The OOD first draw is deliberately reported as a negative result: SFT
is 0.30 points better. Group-RL nevertheless overtakes SFT by $k=8$ and raises
Pass@20 from 35.20\% to 41.70\%. The observed advantage is therefore not a
$k=256$-only effect, but neither is it a uniform improvement in one-shot
generation.

\begin{table}[t]
  \centering
  \caption{Budget-transparent direct de novo generation. Raw strict is the
  candidate-level strict-success fraction; Pass@$k$ is the percentage of
  property programs solved at least once within the first $k$ model draws.}
  \label{tab:budget-scaling}
  \scriptsize
  \setlength{\tabcolsep}{3.5pt}
  \begin{tabular}{lrrrrrr}
    \toprule
    Evaluation & $k$ & \multicolumn{2}{c}{Raw strict (\%)} & \multicolumn{3}{c}{Pass@$k$ (\%)} \\
    \cmidrule{3-4}\cmidrule{5-7}
    & & SFT & Group-RL & SFT & Group-RL & $\Delta$ \\
    \midrule
    2p--7p & 1  & 7.02 & \textbf{8.93} & 7.02 & \textbf{8.93} & +1.91 \\
    2p--7p & 8  & 7.26 & \textbf{9.28} & 35.37 & \textbf{41.72} & +6.35 \\
    2p--7p & 20 & 7.30 & \textbf{9.41} & 54.88 & \textbf{62.52} & +7.64 \\
    \midrule
    OOD & 1  & \textbf{3.50} & 3.20 & \textbf{3.50} & 3.20 & -0.30 \\
    OOD & 8  & 3.20 & \textbf{3.78} & 19.70 & \textbf{23.60} & +3.90 \\
    OOD & 20 & 3.14 & \textbf{3.79} & 35.20 & \textbf{41.70} & +6.50 \\
    \bottomrule
  \end{tabular}
\end{table}

Raw validity explains part, but not all, of this change. Across the matched
2p--7p stream, Group-RL raises candidate validity from 29.04\% to 34.25\% at
$k=20$ while also increasing the fraction of valid molecules that satisfy all
active clauses. OOD validity remains low and nearly unchanged (11.51\% versus
11.60\%), even though OOD Pass@20 improves by 6.50 points. Thus the OOD gain
comes from reallocating probability among the small valid subset, rather than
from solving molecular syntax in general. Full validity and uniqueness
diagnostics are reported in Appendix~\ref{sec:appendix-validity}.

\subsection{The gain persists as programs become more compositional}
\label{sec:results-complexity}

Absolute success decreases with program length for both policies, confirming
that property count defines a genuine compositional difficulty axis.
Table~\ref{tab:complexity-scaling} shows, however, that Group-RL improves every
property count at both practical budgets. At $k=20$, the gain is 4.9 points on
6p programs and 8.4 points on 7p programs; the largest increase, 11.8 points,
occurs at 5p. The aggregate improvement therefore cannot be explained by easy
two-property requests dominating the benchmark.

\begin{table}[t]
  \centering
  \caption{Compositional scaling on the 2p--7p benchmark. Values are the
  percentage of programs solved within the indicated candidate budget.}
  \label{tab:complexity-scaling}
  \small
  \setlength{\tabcolsep}{5pt}
  \begin{tabular}{crrrrrr}
    \toprule
    & \multicolumn{3}{c}{Pass@8} & \multicolumn{3}{c}{Pass@20} \\
    \cmidrule{2-4}\cmidrule{5-7}
    Properties & SFT & Group-RL & $\Delta$ & SFT & Group-RL & $\Delta$ \\
    \midrule
    2 & 58.6 & \textbf{68.1} & +9.5 & 78.6 & \textbf{85.6} & +7.0 \\
    3 & 45.4 & \textbf{52.6} & +7.2 & 66.0 & \textbf{73.7} & +7.7 \\
    4 & 35.7 & \textbf{42.9} & +7.2 & 60.3 & \textbf{66.3} & +6.0 \\
    5 & 30.2 & \textbf{36.1} & +5.9 & 48.7 & \textbf{60.5} & +11.8 \\
    6 & 23.1 & \textbf{27.1} & +4.0 & 42.6 & \textbf{47.5} & +4.9 \\
    7 & 19.2 & \textbf{23.5} & +4.3 & 33.1 & \textbf{41.5} & +8.4 \\
    \bottomrule
  \end{tabular}
\end{table}

The OOD split contains reversed directions, extreme targets, and property
combinations that are rare or absent during training. Group-RL improves the
aggregate OOD stream from $k=8$ onward, but the improvement is not uniform
across OOD strata. In particular, a syntax-safe decoder raises overall OOD
validity and Pass@20 but regresses on the small OOD 7p subset. We therefore
interpret the completed evidence as improved compositional sampling, not as a
claim that high-order extrapolation or validity has been solved.

\subsection{Comparison with matched de novo generation}
\label{sec:results-denovo-external}

Table~\ref{tab:sketchmol-fair40} places the direct MolProgram policy in the
SketchMol multi-property protocol. Each method receives 40 attempts for each
of 1,000 held-out target vectors at every complexity, and a condition succeeds
when at least one valid molecule satisfies every active tolerance. We import
only the published SketchMol row because it uses the matched target and
tolerance contract; other external methods remain blank until rerun rather than
being compared through incompatible datasets. MolProgram is stronger on 2p and
3p and is within 0.9 points at 4p. SketchMol leads from 5p onward and by 5.0
points on the overall average. The external comparison therefore identifies
high-order generation as the primary remaining de novo gap and does not support
an overall state-of-the-art claim.

\begin{table*}[t]
  \centering
  \caption{Matched comparison under the SketchMol multi-property de novo
  protocol. Strict success (\%) uses 1,000 targets per complexity, the same
  property tolerances, and 40 candidates per condition. Dashes denote pending
  matched reruns, not reported failures. Bold marks the best completed result.}
  \label{tab:sketchmol-fair40}
  \scriptsize
  \setlength{\tabcolsep}{3.6pt}
  \begin{tabular}{lllrrrrrrr}
    \toprule
    Method & Molecular representation & Evaluation status & Avg. & 2p & 3p & 4p & 5p & 6p & 7p \\
    \midrule
    SPMM & property vector + SMILES & matched rerun pending
      & -- & -- & -- & -- & -- & -- & -- \\
    STGG+ & spanning-tree sequence & matched rerun pending
      & -- & -- & -- & -- & -- & -- & -- \\
    Graph DiT & molecular graph & matched rerun pending
      & -- & -- & -- & -- & -- & -- & -- \\
    SketchMol & molecular image & published matched reference
      & \textbf{73.1} & 80.4 & 76.8 & \textbf{73.6} & \textbf{71.6} & \textbf{67.8} & \textbf{68.5} \\
    \midrule
    MolProgram (direct Group-RL) & direct SMILES & ours, matched rerun
      & 68.1 & \textbf{91.7} & \textbf{83.2} & 72.7 & 63.6 & 51.0 & 46.5 \\
    \bottomrule
  \end{tabular}
\end{table*}

\subsection{One transition language covers empty and source graphs}
\label{sec:results-transition-audit}

Before measuring task performance, we audit whether one typed transition
language can faithfully express trajectories for both initial-state regimes.
Table~\ref{tab:transition-audit} uses the frozen mixed training pool and applies
the same serializer and deterministic interpreter to every trajectory. The
language covers 96.5\% of all rows under the fixed 188-token limit. Editing
coverage is 100\%, and 99.8\% of covered editing programs reconstruct their
selected executable targets exactly after canonicalization. De novo coverage
is lower because atom-by-atom construction produces longer programs; exact
reconstruction is additionally affected by stereochemical details that are not
fully represented in the current bond state. These are representation
diagnostics, not held-out molecular-design accuracies.

\begin{table}[t]
  \centering
  \caption{Unified molecular-transition representation audit. ``Covered''
  denotes a serializable program within the fixed token limit; exact
  reconstruction is measured only among covered rows.}
  \label{tab:transition-audit}
  \small
  \setlength{\tabcolsep}{5pt}
  \begin{tabular}{lrrrr}
    \toprule
    Initial graph & Input rows & Covered rows & Coverage (\%) & Exact recon. (\%) \\
    \midrule
    Empty (de novo) & 288 & 261 & 90.6 & 80.5 \\
    Source (editing) & 480 & 480 & 100.0 & \textbf{99.8} \\
    \midrule
    Overall & 768 & 741 & \textbf{96.5} & 93.0 \\
    \bottomrule
  \end{tabular}
\end{table}

The audit isolates the intended meaning of \emph{unified}: the initial graph
changes, but the property-program interface, initial-state observation,
autoregressive decoder, checkpoint, action vocabulary, and interpreter do not.
The initial-state token is computed from graph occupancy and does not select a
task-specific head. The first held-out unified-policy gate evaluates raw
candidate streams on hard 6p/7p de novo conditions and on all ten
MolEdit-Instruct task families. Its result is asymmetric. Editing obtains at
least one valid molecule for every input and reaches 26.5\%, 47.5\%, and 54.0\%
macro strict accuracy at $k=1,8,20$, respectively. Across the complete unranked
20-candidate pool, validity is 65.55\% and strict candidate-level accuracy is
17.75\%. In contrast, only 3 of 1,280 hard de novo action sequences execute to
valid molecules (0.234\%), and none satisfies all six or seven properties. The
shared policy therefore learns short source-conditioned transitions but not
long-horizon empty-graph construction.

\begin{table}[t]
  \centering
  \caption{First single-checkpoint unified-policy gate. Candidate validity is
  measured over all unselected draws; Strict@1 and Pass@$k$ are averaged over
  conditions (and then over task families for editing).}
  \label{tab:unified-policy-gate}
  \small
  \setlength{\tabcolsep}{4pt}
  \begin{tabular}{lrrrrr}
    \toprule
    Evaluation & Conditions & Cand. validity & Strict@1 & Pass@8 & Pass@20 \\
    \midrule
    Hard de novo (6p/7p) & 64 & 0.23 & 0.0 & 0.0 & 0.0 \\
    MolEdit Table 1 pilot & 200 & 65.55 & \textbf{26.5} & \textbf{47.5} & \textbf{54.0} \\
    \bottomrule
  \end{tabular}
\end{table}

\subsection{Source-conditioned editing against published baselines}
\label{sec:results-editing}

Table~\ref{tab:moledit-full} retains the complete MolEdit-Instruct comparison:
ten tasks, validity, and overall/valid-conditional accuracy at source
similarity thresholds 0.65 and 0.15. Every published method is evaluated from
20 generated molecules per input and candidate-wise over the unranked pool.
This is intentionally different from Pass@$k$ or a property-selected output.
The unified MolProgram row reports the complete P6 candidate pool on a bounded
20-input-per-task pilot and is marked as preliminary rather than treated as a
full-scale matched result. We omit FCD because it is not part of our evaluation
contract.

\begin{table*}[h]
  \centering
  \caption{Strict
  $\operatorname{AnyAcc}_{\rm all}^{@K}(.65)$ by MolEdit-Instruct task
  ($n=10$ each). Bold marks the better checkpoint at a fixed budget; ties are
  left unbolded. HBA and RB correspond to the benchmark's Haccept and Rotbonds
  controls.}
  \label{tab:editing-by-task}
  \scriptsize
  \setlength{\tabcolsep}{4pt}
  \begin{tabular}{lrrrrrr}
    \toprule
    Task & P17 @1 & P18 @1 & P17 @4 & P18 @4 & P17 @8 & P18 @8 \\
    \midrule
    DRD2$\downarrow$ + MW$\downarrow$ + SA$\downarrow$
      & \textbf{0.4} & 0.1 & \textbf{0.4} & 0.2 & 0.4 & 0.4 \\
    GSK3B$\uparrow$
      & 0.0 & 0.0 & 0.0 & 0.0 & 0.0 & 0.0 \\
    HBA$\downarrow$ + LogP$\uparrow$
      & 0.2 & 0.2 & 0.2 & \textbf{0.3} & 0.2 & \textbf{0.4} \\
    HBA$\downarrow$ + MW$\downarrow$
      & 0.2 & 0.2 & 0.2 & \textbf{0.3} & 0.2 & \textbf{0.3} \\
    HBA$\downarrow$ + SA$\downarrow$
      & 0.1 & \textbf{0.3} & 0.2 & \textbf{0.6} & 0.3 & \textbf{0.6} \\
    HBA$\uparrow$ + MW$\uparrow$ + QED$\downarrow$
      & 0.2 & 0.2 & 0.2 & 0.2 & 0.3 & 0.3 \\
    MW$\uparrow$
      & 0.3 & \textbf{0.4} & 0.4 & \textbf{0.7} & 0.4 & \textbf{0.9} \\
    QED$\uparrow$ + SA$\downarrow$
      & 0.2 & 0.2 & 0.2 & \textbf{0.3} & 0.2 & \textbf{0.3} \\
    RB$\downarrow$
      & 0.0 & 0.0 & 0.0 & 0.0 & 0.0 & \textbf{0.1} \\
    SA$\downarrow$
      & 0.1 & \textbf{0.2} & 0.1 & \textbf{0.3} & 0.1 & \textbf{0.3} \\
    \midrule
    Macro
      & 0.17 & \textbf{0.18} & 0.19 & \textbf{0.29}
      & 0.21 & \textbf{0.36} \\
    \bottomrule
  \end{tabular}
\end{table*}


The earlier direct-SMILES transfer experiments explain why the unified policy
uses executable transitions rather than a second free-form string mode. A
source-aware direct decoder produces valid molecules but fails to preserve the
input at the strict similarity threshold, whereas a standalone executable
action policy reaches 100\% validity and 26.5\% raw first-candidate strict
accuracy on the paired repair subset. A source-only finalizer can raise the
latter number to 66.0\% at $k=8$, but this is assisted selection and is kept in
the appendix. The P6 row is populated only by raw samples from the single
shared transition checkpoint; its 17.75\% macro strict candidate accuracy is
below MolEditRL overall, although it exceeds MolEditRL on the MW-increase task
at the strict 0.65 threshold.

\subsection{MuMOInstruct IND and OOD task families}
\label{sec:results-mumo}

Tables~\ref{tab:mumo-ind-full} and~\ref{tab:mumo-ood-full} preserve the five IND
and five OOD MuMOInstruct task families and their three complementary metrics:
success rate (SR), similarity among successful outputs (Sim), and relative
improvement among successes (RI) \citep{dey2025gellm3o}. This fuller layout is
necessary because a high property SR can coexist with very low source
similarity. The unified MolProgram cells remain blank until all official
oracles and source-fidelity metrics are exported from the same 20-candidate
run.

\begin{table*}[t]
  \centering
  \caption{MuMOInstruct performance on in-distribution (IND) tasks with 20
  candidates per input. External values are transcribed from Table~3 of
  \citet{dey2025gellm3o}. SR is success rate (\%); Sim and RI are benchmark
  similarity and relative-improvement metrics over successful cases. Dashes
  denote an inapplicable task or a pending matched evaluation. Bold marks the
  best completed value for each task and metric.}
  \label{tab:mumo-ind-full}
  \setlength{\tabcolsep}{1.6pt}
  \renewcommand{\arraystretch}{0.86}
  \resizebox{\textwidth}{!}{%
  \begin{tabular}{lrrrrrrrrrrrrrrr}
    \toprule
    & \multicolumn{3}{c}{BDP} & \multicolumn{3}{c}{BDQ}
    & \multicolumn{3}{c}{BPQ} & \multicolumn{3}{c}{DPQ}
    & \multicolumn{3}{c}{BDPQ} \\
    \cmidrule(lr){2-4}\cmidrule(lr){5-7}\cmidrule(lr){8-10}
    \cmidrule(lr){11-13}\cmidrule(lr){14-16}
    Model & SR$\uparrow$ & Sim$\uparrow$ & RI$\uparrow$
    & SR$\uparrow$ & Sim$\uparrow$ & RI$\uparrow$
    & SR$\uparrow$ & Sim$\uparrow$ & RI$\uparrow$
    & SR$\uparrow$ & Sim$\uparrow$ & RI$\uparrow$
    & SR$\uparrow$ & Sim$\uparrow$ & RI$\uparrow$ \\
    \midrule
    \multicolumn{16}{c}{\textit{General-purpose LLMs}} \\
    Mistral (0-shot) & 6.60 & \textbf{.81} & .68 & 3.00 & \textbf{.76} & .53 & 15.80 & .73 & .51 & 2.20 & \textbf{.65} & .41 & 3.20 & .77 & .87 \\
    Llama (0-shot) & 22.00 & .73 & .74 & 2.20 & .64 & .53 & 28.40 & .64 & .72 & 2.60 & .62 & .32 & 5.20 & \textbf{.80} & .62 \\
    Claude-3.5 (0-shot) & 19.60 & .66 & 1.05 & 13.00 & .62 & 1.14 & 56.00 & .62 & .86 & 11.00 & .54 & .51 & 8.00 & .60 & 1.34 \\
    GPT-4o (0-shot) & 7.80 & .69 & .90 & 2.00 & .69 & .62 & 36.40 & .73 & .42 & 2.80 & .57 & .50 & 1.80 & .71 & .39 \\
    Mistral (5-shot) & 35.20 & .64 & 2.10 & 17.00 & .60 & 2.32 & 68.60 & .63 & .79 & 10.40 & .54 & 1.10 & 11.00 & .69 & .96 \\
    Llama (5-shot) & 35.40 & .57 & 2.71 & 16.60 & .43 & \textbf{5.70} & 34.60 & .70 & .64 & 8.20 & .44 & 3.02 & 9.60 & .54 & 3.45 \\
    Claude-3.5 (5-shot) & 35.40 & .50 & 2.43 & 29.40 & .43 & 3.80 & 76.80 & .53 & 1.24 & 29.20 & .37 & 2.87 & 20.80 & .35 & 3.53 \\
    GPT-4o (1-shot) & 9.40 & .69 & .79 & 7.60 & .66 & .61 & 40.00 & \textbf{.75} & .41 & 7.00 & .62 & .44 & 3.40 & .70 & .61 \\
    \midrule
    \multicolumn{16}{c}{\textit{Foundation LLMs for chemistry}} \\
    ChemLLM & .20 & .17 & 1.20 & 1.00 & .55 & .82 & 4.80 & .29 & .96 & .60 & .28 & .42 & .00 & n/a & n/a \\
    LlaSMol$_{\mathrm{Mistral}}$ & 43.60 & .62 & 1.09 & 31.40 & .66 & .93 & 86.00 & .58 & .84 & 24.00 & .57 & .61 & 14.00 & .62 & 1.03 \\
    \midrule
    \multicolumn{16}{c}{\textit{Task-specific non-LLM}} \\
    Prompt-MolOpt & 12.20 & .12 & \textbf{7.46} & 23.20 & .10 & 5.40 & 15.80 & .10 & 1.50 & 23.60 & .10 & \textbf{5.46} & 6.60 & .11 & \textbf{5.36} \\
    \midrule
    \multicolumn{16}{c}{\textit{Task-specific LLMs}} \\
    GeLLM$^3$O-3$_{\mathrm{Mistral}}$ & 84.80 & .47 & 4.30 & 87.00 & .47 & 5.61 & 93.00 & .46 & 1.49 & 62.80 & .37 & 3.87 & -- & -- & -- \\
    GeLLM$^3$O-3$_{\mathrm{Llama}}$ & \textbf{86.80} & .48 & 4.38 & \textbf{90.00} & .46 & 5.66 & 94.00 & .50 & 1.38 & 60.60 & .44 & 3.76 & -- & -- & -- \\
    GeLLM$^3$O-4$_{\mathrm{Mistral}}$ & 71.60 & .49 & 3.27 & 57.40 & .55 & 2.56 & 90.20 & .46 & 1.41 & 54.00 & .44 & 3.02 & 30.00 & .48 & 3.44 \\
    GeLLM$^3$O-4$_{\mathrm{Llama}}$ & 53.60 & .63 & 1.94 & 48.60 & .59 & 1.29 & 93.40 & .59 & 1.12 & 39.60 & .57 & 1.32 & 28.00 & .66 & 1.02 \\
    \midrule
    \multicolumn{16}{c}{\textit{Generalist LLMs}} \\
    GeLLM$^3$O-P(3)$_{\mathrm{Mistral}}$ & 75.60 & .56 & 3.31 & 79.40 & .53 & 4.52 & 93.20 & .55 & 1.23 & 57.20 & .50 & 2.22 & -- & -- & -- \\
    GeLLM$^3$O-P(3)$_{\mathrm{Llama}}$ & 77.40 & .51 & 3.16 & 76.40 & .57 & 4.41 & 95.40 & .50 & 1.46 & 63.40 & .49 & 2.46 & -- & -- & -- \\
    GeLLM$^3$O-P(4)$_{\mathrm{Mistral}}$ & 81.40 & .55 & 3.95 & 82.60 & .56 & 5.24 & 96.20 & .52 & 1.52 & \textbf{66.60} & .53 & 2.41 & \textbf{57.40} & .52 & 3.04 \\
    GeLLM$^3$O-P(4)$_{\mathrm{Llama}}$ & 80.40 & .54 & 3.60 & 81.40 & .56 & 4.81 & 93.80 & .47 & 1.64 & 61.40 & .50 & 2.02 & 49.80 & .48 & 3.26 \\
    GeLLM$^3$O-P(6)$_{\mathrm{Mistral}}$ & 83.00 & .57 & 3.60 & 85.80 & .59 & 4.78 & \textbf{96.80} & .53 & 1.48 & 60.80 & .54 & 2.16 & 54.00 & .54 & 3.09 \\
    GeLLM$^3$O-P(6)$_{\mathrm{Llama}}$ & 77.00 & .53 & 3.73 & 79.60 & .56 & 5.05 & 95.00 & .47 & \textbf{1.66} & 57.00 & .49 & 2.50 & 52.20 & .49 & 3.48 \\
    \midrule
    \multicolumn{16}{c}{\textit{Ours}} \\
    MolProgram (unified) & -- & -- & -- & -- & -- & -- & -- & -- & -- & -- & -- & -- & -- & -- & -- \\
    \bottomrule
  \end{tabular}}
\end{table*}

\begin{table*}[t]
  \centering
  \caption{MuMOInstruct performance on out-of-distribution (OOD) tasks with
  20 candidates per input. External values are transcribed from Table~4 of
  \citet{dey2025gellm3o}; notation follows
  Table~\ref{tab:mumo-ind-full}.}
  \label{tab:mumo-ood-full}
  \setlength{\tabcolsep}{1.6pt}
  \renewcommand{\arraystretch}{0.90}
  \resizebox{\textwidth}{!}{%
  \begin{tabular}{lrrrrrrrrrrrrrrr}
    \toprule
    & \multicolumn{3}{c}{MPQ} & \multicolumn{3}{c}{BDMQ}
    & \multicolumn{3}{c}{BHMQ} & \multicolumn{3}{c}{BMPQ}
    & \multicolumn{3}{c}{HMPQ} \\
    \cmidrule(lr){2-4}\cmidrule(lr){5-7}\cmidrule(lr){8-10}
    \cmidrule(lr){11-13}\cmidrule(lr){14-16}
    Model & SR$\uparrow$ & Sim$\uparrow$ & RI$\uparrow$
    & SR$\uparrow$ & Sim$\uparrow$ & RI$\uparrow$
    & SR$\uparrow$ & Sim$\uparrow$ & RI$\uparrow$
    & SR$\uparrow$ & Sim$\uparrow$ & RI$\uparrow$
    & SR$\uparrow$ & Sim$\uparrow$ & RI$\uparrow$ \\
    \midrule
    \multicolumn{16}{c}{\textit{General-purpose LLMs}} \\
    Mistral (0-shot) & 11.20 & .57 & .48 & 1.20 & .68 & .37 & 12.71 & .73 & 1.90 & 12.57 & .61 & .54 & 21.88 & \textbf{.72} & .72 \\
    Llama (0-shot) & 25.80 & .44 & .61 & 1.20 & \textbf{.76} & .30 & 11.02 & .74 & .68 & 16.75 & .51 & .57 & 15.62 & .47 & .60 \\
    Claude-3.5 (0-shot) & 17.40 & .49 & .52 & 15.00 & .57 & .87 & 38.98 & .51 & 2.35 & 44.50 & .55 & .85 & 38.54 & .54 & 1.01 \\
    GPT-4o (0-shot) & 19.40 & \textbf{.61} & .35 & 1.60 & .67 & .18 & 17.80 & .66 & .65 & 25.13 & .67 & .35 & 20.83 & .67 & .37 \\
    Mistral (5-shot) & 59.60 & .54 & .57 & 20.40 & .59 & 1.65 & 34.75 & .70 & 1.31 & 49.21 & .62 & .73 & 46.88 & .66 & .91 \\
    Llama (5-shot) & 34.80 & .57 & .53 & 16.80 & .39 & 3.22 & 36.44 & .67 & 1.13 & 31.94 & .66 & .60 & 33.33 & .68 & .61 \\
    Claude-3.5 (5-shot) & 50.60 & .49 & .71 & 30.40 & .49 & 2.32 & 52.54 & .48 & 2.52 & 52.36 & .46 & 1.08 & 65.62 & .48 & 1.32 \\
    GPT-4o (1-shot) & 21.40 & .60 & .48 & 6.20 & .70 & .75 & 14.41 & \textbf{.75} & .67 & 24.08 & \textbf{.68} & .45 & 25.00 & .70 & .54 \\
    \midrule
    \multicolumn{16}{c}{\textit{Foundation LLMs for chemistry}} \\
    ChemLLM & 6.20 & .31 & .61 & .00 & n/a & n/a & 1.69 & .53 & .74 & 5.24 & .24 & .72 & 3.12 & .58 & .51 \\
    LlaSMol$_{\mathrm{Mistral}}$ & 76.40 & .55 & .53 & 28.20 & .66 & .52 & 53.39 & .62 & 1.14 & 64.92 & .58 & .57 & 53.12 & .62 & .70 \\
    \midrule
    \multicolumn{16}{c}{\textit{Generalist LLMs}} \\
    GeLLM$^3$O-P(6)$_{\mathrm{Mistral}}$ & \textbf{95.20} & .53 & .85 & \textbf{79.00} & .56 & 3.10 & 86.44 & .54 & 2.58 & 91.10 & .53 & 1.06 & 91.67 & .55 & 1.42 \\
    GeLLM$^3$O-P(6)$_{\mathrm{Llama}}$ & 93.60 & .48 & \textbf{.91} & 74.20 & .55 & \textbf{3.25} & \textbf{93.22} & .49 & \textbf{3.57} & \textbf{95.29} & .49 & \textbf{1.20} & \textbf{97.92} & .46 & \textbf{1.76} \\
    \midrule
    \multicolumn{16}{c}{\textit{Ours}} \\
    MolProgram (unified) & -- & -- & -- & -- & -- & -- & -- & -- & -- & -- & -- & -- & -- & -- & -- \\
    \bottomrule
  \end{tabular}}
  \vspace{2pt}
  \parbox{\textwidth}{\scriptsize Task codes follow MuMOInstruct:
  B=BBBP, D=DRD2, H=HIA, M=mutagenicity, P=penalized LogP, and Q=QED.}
\end{table*}


Taken together, the completed results support two bounded conclusions.
Group-relative post-training increases the probability that a modest direct
sample set solves a compositional property program, including at 6p, 7p, and
on aggregate OOD requests. Separately, one typed graph-transition language can
represent both empty-graph construction and source-graph editing with high
coverage and near-exact editing execution. The first jointly trained checkpoint
does not retain both advantages simultaneously: it provides non-trivial editing
but collapses on long de novo action programs. Unified representation is
therefore established, while unified policy optimization remains an open
problem rather than a completed performance claim.


\section{Discussion and Limitations}
% To be written.

\section{Conclusion}
% To be written.

\appendix
\section{Additional Results and Diagnostics}
\label{sec:appendix-results}

\subsection{Raw validity across candidate budgets}
\label{sec:appendix-validity}

Table~\ref{tab:raw-validity} reports raw candidate validity separately from
strict property success. Group-RL improves in-distribution validity by roughly
five points, but OOD validity remains close to 12\% for both policies. Nearly
every valid draw is unique, indicating that invalid SMILES syntax rather than
conventional mode collapse is the primary sample-efficiency bottleneck.

\begin{table}[h]
  \centering
  \caption{Raw candidate validity (\%) at matched generation prefixes.}
  \label{tab:raw-validity}
  \small
  \begin{tabular}{llrrr}
    \toprule
    Evaluation & Method & $k=1$ & $k=8$ & $k=20$ \\
    \midrule
    2p--7p & SFT & 29.20 & 28.85 & 29.04 \\
    2p--7p & Group-RL & \textbf{33.53} & \textbf{33.86} & \textbf{34.25} \\
    OOD & SFT & \textbf{12.50} & \textbf{11.92} & 11.51 \\
    OOD & Group-RL & 12.10 & 11.63 & \textbf{11.60} \\
    \bottomrule
  \end{tabular}
\end{table}

\subsection{Larger rollout groups are not sufficient}
\label{sec:appendix-g64}

Table~\ref{tab:g64-ablation} compares the main group-16 policy with a
hard-only continuation using 64 rollouts per program on a fixed 6p/7p subset.
The larger group improves validity but reduces Pass@20 from 50.78\% to 45.31\%,
driven by a large 7p regression. We therefore retain group 16 and report group
64 as a negative ablation rather than selecting it on raw validity.

\begin{table}[h]
  \centering
  \caption{Hard-program rollout-group ablation (\%). G64 continues training
  from the G16 checkpoint using only 6p/7p programs.}
  \label{tab:g64-ablation}
  \scriptsize
  \setlength{\tabcolsep}{3.5pt}
  \begin{tabular}{lrrrrrr}
    \toprule
    Method & $k$ & Raw strict & Pass@$k$ & 6p Pass@$k$ & 7p Pass@$k$ & Validity \\
    \midrule
    SFT & 8 & 3.47 & 21.88 & 21.88 & 21.88 & 28.27 \\
    Group-RL G16 & 8 & \textbf{4.74} & \textbf{28.12} & 31.25 & \textbf{25.00} & 32.91 \\
    Group-RL G64 & 8 & 4.44 & 25.78 & \textbf{32.81} & 18.75 & \textbf{36.38} \\
    \midrule
    SFT & 20 & 3.63 & 37.11 & 39.84 & 34.38 & 29.00 \\
    Group-RL G16 & 20 & 4.75 & \textbf{50.78} & \textbf{50.78} & \textbf{50.78} & 33.26 \\
    Group-RL G64 & 20 & \textbf{5.02} & 45.31 & \textbf{50.78} & 39.84 & \textbf{35.63} \\
    \bottomrule
  \end{tabular}
\end{table}

\subsection{MolEdit-Instruct comparison}
\label{sec:appendix-moledit}

Table~\ref{tab:moledit-horizontal} places MolProgram alongside seven published
MolEdit-Instruct baselines. External values are ten-task macro-averages
recomputed from Table~1 of \citet{zhuang2025moleditrl}. The comparison exposes
the boundary of direct decoding: MolProgram's first raw candidate is much less
valid and less source-preserving than every specialized editing model. A
property-aware 20-candidate selector improves relaxed accuracy but remains
assisted inference and does not repair the strict 0.65 threshold.

\begin{table}[h]
  \centering
  \caption{MolEdit-Instruct ten-task macro-average (\%).
  $\operatorname{Acc}_{\mathrm{all}}(t)$ requires both the requested property
  directions and source Tanimoto similarity of at least $t$.}
  \label{tab:moledit-horizontal}
  \scriptsize
  \setlength{\tabcolsep}{4pt}
  \begin{tabular}{llrrr}
    \toprule
    Method & Protocol & Validity & $\operatorname{Acc}_{\mathrm{all}}(.65)$ & $\operatorname{Acc}_{\mathrm{all}}(.15)$ \\
    \midrule
    BioT5 & reported & \textbf{100.0} & 0.0 & 29.2 \\
    MolGen & reported & \textbf{100.0} & 3.3 & 22.2 \\
    REINVENT4 & reported & 64.1 & 12.5 & 25.0 \\
    DrugAssist & reported & 97.9 & 34.9 & 41.1 \\
    GeLLM3O$_{\mathrm{Mistral}}$ & reported & 90.2 & 15.2 & 35.4 \\
    GeLLM3O$_{\mathrm{Llama}}$ & reported & 90.1 & 14.5 & 40.4 \\
    MolEditRL & reported & 96.6 & \textbf{45.0} & \textbf{72.7} \\
    \midrule
    MolProgram (Group-RL) & first raw candidate & 23.4 & 0.0 & 3.9 \\
    MolProgram (SFT) & assisted, $k=20$ & 96.4 & 0.0 & 17.9 \\
    MolProgram (Group-RL) & assisted, $k=20$ & 97.2 & 0.0 & 19.0 \\
    \bottomrule
  \end{tabular}
\end{table}

\subsection{MuMO source-fidelity diagnostic}
\label{sec:appendix-mumo}

Under the official 20-candidate MuMO oracle, MolProgram reaches 50.4\% overall
property SR, with 32.3\% IND and 68.8\% OOD SR. However, mean source
similarity among successful outputs is only 0.098, and no input passes the
additional 0.4 source-similarity diagnostic. The policy therefore often finds
a molecule with the desired properties by changing the molecular identity
rather than editing the input locally. MuMO SR is useful evidence for property
control, but it must not be presented as source-preserving editing success.

\begin{table}[h]
  \centering
  \caption{MolProgram on MuMOInstruct under the official property SR and an
  additional source-fidelity diagnostic.}
  \label{tab:mumo-fidelity}
  \small
  \begin{tabular}{lrrrrr}
    \toprule
    Method & Overall SR & IND SR & OOD SR & Sim(success) & Sim$\geq .4$ \\
    \midrule
    MolProgram (Group-RL) & 50.4 & 32.3 & 68.8 & 0.098 & 0.0 \\
    \bottomrule
  \end{tabular}
\end{table}

\subsection{Executable actions improve editing reachability}
\label{sec:appendix-actions}

A source-similarity reward applied to free-form direct SMILES Group-RL did not
produce any candidates above the required source-similarity threshold. In a
paired 200-example validation experiment, replacing free-form editing with
executable GraphEditDSL actions raises raw $k=1$ validity from 31.0\% to
100\% and strict $\operatorname{Acc}_{\mathrm{all}}(.65)$ from 0\% to 20.0\%,
without reducing held-out de novo strict success. At $k=8$, an explicitly
assisted finalizer reaches 43.5\% strict accuracy. These validation results
motivate source-consistent action policies, but are not part of the main
MolProgram claim.

\subsection{Frozen slots for ongoing repair experiments}
\label{sec:appendix-pending}

We retain empty cells for the preregistered repair runs rather than filling
them with selected-molecule validity or results from a different evaluation
subset.

\begin{table}[h]
  \centering
  \caption{Syntax-safe decoding repair. Pending rows remain blank until the
  paired P2 evaluation completes.}
  \label{tab:validity-repair}
  \small
  \begin{tabular}{llrrr}
    \toprule
    Evaluation & Decoder & Raw validity@1 & Raw strict@1 & Pass@8 \\
    \midrule
    2p--7p & legacy Group-RL & 33.53 & 8.93 & 41.72 \\
    2p--7p & syntax-safe Group-RL & -- & -- & -- \\
    OOD & legacy Group-RL & 12.10 & 3.20 & 23.60 \\
    OOD & syntax-safe Group-RL & -- & -- & -- \\
    \bottomrule
  \end{tabular}
\end{table}

\begin{table}[h]
  \centering
  \caption{Source-consistent executable editing. Pending rows use source-only
  consistency ranking; target molecules and output property oracles do not
  enter raw selection.}
  \label{tab:edit-repair}
  \scriptsize
  \setlength{\tabcolsep}{3pt}
  \begin{tabular}{lrrr}
    \toprule
    Edit policy & Validity@1 & $\operatorname{Acc}_{\mathrm{all}}(.65)$ & $\operatorname{Acc}_{\mathrm{all}}(.15)$ \\
    \midrule
    Protected action policy & 100.0 & 20.0 & 20.0 \\
    Source-consistent action policy & -- & -- & -- \\
    Strong source-consistent action policy & -- & -- & -- \\
    \bottomrule
  \end{tabular}
\end{table}


\bibliography{references}
\bibliographystyle{iclr2027_conference}

\end{document}
